%% kapa: exam=1 round=24 year=2013 period=1 form=A subject=economics subject_ko=경제원론 seq=4 number=44
%% source: Q-Net 감정평가사 기출문제 (cst00309) artlSeq=5073802 file=extracted/1008/5073802/2013년 제24회 감정평가사 1차시험 기출문제/1교시(A형).hwp
%% conversion: kordoc 4.24.51 -> tools/gen1.py (auto); review=proofread-sonnet; flags=none
\begin{question}{44}
효용극대화를 추구하는 소비자 甲의 효용함수는 $U(x , y)=\min\{2x , y\}$이다. 甲의 수요에 관한 설명으로 옳지 않은 것은? (단, 甲은 $X$재와 $Y$재만 소비하고, $x$는 $X$재 소비량, $y$는 $Y$재 소비량이다.)
\begin{choices}
  \item 상대가격변화에 따른 대체효과는 0이다.
  \item $X$재 수요의 소득탄력성은 1이다.
  \item 가격소비곡선은 수직선의 형태를 갖는다.
  \item 소득소비곡선은 원점에서 출발하는 직선의 형태를 갖는다.
  \item 甲의 효용함수는 1차 동차함수이다.
\end{choices}
\end{question}
\answer{3}
